% TOP_PROBLEM: {"problemId": "problem.corona-problem-for-the-unit-polydisc", "canonicalTitle": "Corona problem for the unit polydisc", "releaseRank": 349, "primaryDomain": "analysis_pde", "primaryDomainLabel": "Analysis & PDE", "displayStatus": "open_with_solved_subcases", "releaseStatus": "open_with_solved_subcases"}
% !TeX program = lualatex
% ATTEMPT_STATUS: unresolved
% Model: GPT-6 Astra Ultra; continuation dated 27 September 2026.
\documentclass[11pt]{article}
\usepackage[margin=25mm]{geometry}
\usepackage{fontspec}
\setmainfont{Latin Modern Roman}
\usepackage{amsmath,amssymb,mathtools}
\usepackage{unicode-math}
\setmathfont{Latin Modern Math}
\usepackage{xurl,hyperref}
\hypersetup{colorlinks=true,urlcolor=blue}
\setcounter{secnumdepth}{0}
\setlength{\parindent}{0pt}
\setlength{\parskip}{5pt}
\setlength{\emergencystretch}{3em}
\begin{document}
\section*{\texorpdfstring{Corona problem for the unit polydisc}{Corona problem for the unit polydisc}}\label{349}
\subsection{Definitions and mathematical statement}
Let $\mathbb D^n=\{z\in{\mathbb{C}}^n:|z_k|<1\text{ for all }k\}$ and let $H^\infty(\mathbb D^n)$ be the Banach algebra of bounded holomorphic scalar functions with the supremum norm. For every $n\ge2$, finite $m\ge1$, and $f_1,\ldots,f_m$ in this algebra satisfying
\[
 \inf_{z\in\mathbb D^n}\sum_{j=1}^m|f_j(z)|>0,
\]
the corona problem asks for $g_1,\ldots,g_m\in H^\infty(\mathbb D^n)$ such that $\sum_jf_jg_j=1$. The maximal ideal space consists of nonzero multiplicative linear functionals with its weak-star topology; the equivalent density statement says that point evaluations are dense there. No prescribed norm bound on the $g_j$, boundary continuity, matrix-valued extension, or stronger operator inequality is included.
\par\smallskip\textit{Formulation and concept references:} \href{https://arxiv.org/pdf/2512.23733v2}{[S1]}.

\subsection{Short English statement}
Can bounded holomorphic functions on every higher-dimensional polydisc that stay uniformly away from a common zero always solve a bounded holomorphic B\'ezout equation?
\subsection{Research attempt}
\paragraph{Scope and the source's extra hypothesis.}
The arbitrary-data problem is \textbf{UNRESOLVED} here. The primary
source [S1], version 2 of 3 January 2026, explicitly distinguishes
the unsolved multidimensional corona problem from its gluing theorem.
That theorem assumes local solvability on a special cover of a maximal
ideal space. Local solvability at ordinary points of the polydisc is
not the same hypothesis. We prove a substantial restricted algebra
case, exhibit why it does not exhaust all bounded holomorphic data,
and isolate a uniform estimate sufficient to pass beyond it.

Write $A=H^\infty(\mathbb D^n)$ and suppose
$\sum_j|f_j|\geq\delta>0$. If $m=1$, the reciprocal $1/f_1$ is
holomorphic and bounded by $1/\delta$, giving a complete solution.
More generally, if a constant linear combination $F=\sum c_jf_j$
is bounded away from zero, $g_j=c_j/F$ solves the equation. The
general corona hypothesis does not provide a single such combination.

\paragraph{The maximal-ideal reformulation, with its quantifiers.}
We use the elementary Banach-algebra theorem that characters of a
complex commutative unital Banach algebra are continuous of norm one,
and every proper ideal lies in the kernel of a character. Here is
the relevant reason the theorem applies even to a nonclosed ideal.
Its closure cannot contain the unit: an element within distance less
than one of the unit is invertible by the Neumann series and would
make the ideal the whole algebra. A maximal ideal containing it is
therefore closed; its Banach division quotient is $\mathbb C$ by the
Gelfand--Mazur theorem.

If ordinary point evaluations are dense in $M(A)$, continuity of
$\chi\mapsto\sum_j|\chi(f_j)|$ extends the lower bound $\delta$
to every character. The ideal generated by the $f_j$ then cannot
be proper. Thus there are coefficients in $A$ with $\sum f_jg_j=1$.
Conversely, if a character $\chi$ were outside the closure of point
evaluations, some basic neighborhood, specified by finitely many
functions $h_j$, would miss all evaluations. Consequently
\[
 \inf_z\max_j|h_j(z)-\chi(h_j)|>0.
\]
A corona solution for $f_j=h_j-\chi(h_j)$ would give $0=1$ after
applying $\chi$. This proves the stated equivalence. It does not
say that every character is itself evaluation at an interior point.

\paragraph{A complete theorem for the uniform slice algebra.}
Let $\mathcal S_n$ be the closure in $A$ of finite sums of products
\[
 \prod_{k=1}^n a_k(z_k),\qquad a_k\in H^\infty(\mathbb D).
 \tag{349.1}
\]
This is a closed unital subalgebra. We use the one-dimensional corona
theorem, equivalently density of disk evaluations in
$M(H^\infty(\mathbb D))$, as a declared external theorem. The
following deduction from it is supplied in detail.

Let $\chi\in M(\mathcal S_n)$. Restricting $\chi$ to the copy of
$H^\infty(\mathbb D)$ in the $k$th coordinate gives a character
$\chi_k$. For a product in (349.1), multiplicativity gives
$\chi(\prod a_k)=\prod\chi_k(a_k)$. Given finitely many elements
$F_1,\ldots,F_b$ of $\mathcal S_n$ and $\varepsilon>0$, approximate
each uniformly by a finite sum of products. Only finitely many
one-variable factors occur. In coordinate $k$, disk corona density
provides $z_k\in\mathbb D$ approximating $\chi_k$ simultaneously
on all those factors, to any prescribed tolerance.

For bounded factors, the elementary telescoping product identity
\[
 \prod_{k=1}^n a_k-\prod_{k=1}^n b_k
 =\sum_{k=1}^n(a_k-b_k)
       \prod_{j<k}b_j\prod_{j>k}a_j
 \tag{349.2}
\]
controls the resulting error in each product. Taking the finitely
many tolerances sufficiently small, and using $\|\chi\|=1$ on
the initial uniform approximation errors, yields
$|F_i(z)-\chi(F_i)|<\varepsilon$ for every $i$. Evaluations on
$\mathbb D^n$ are therefore dense in $M(\mathcal S_n)$.

Apply the preceding maximal-ideal argument within $\mathcal S_n$:
if all the corona data lie in this algebra, the coefficients can
also be chosen in it. This includes arbitrary finite separated sums
and their uniform limits. Boundary continuity is not required for
the one-variable factors. The imported deep step is exactly the
one-dimensional corona theorem, not a multidimensional assertion.

For comparison, let the generators be grouped by coordinates,
$f_{k,j}=a_{k,j}(z_k)$. Then
\[
 \inf_{z\in\mathbb D^n}\sum_{k,j}|a_{k,j}(z_k)|
 =\sum_k\inf_{w\in\mathbb D}\sum_j|a_{k,j}(w)|.
\]
The inequality in one direction is immediate; for the other choose
each coordinate separately within an arbitrary error of its infimum.
A positive left side forces a positive one-variable corona lower
bound in at least one group. Solving in that group and setting all
other coefficients to zero gives a simpler explicit reduction.

\paragraph{The slice algebra is genuinely smaller than the full algebra.}
For $F\in\mathcal S_2$, the family of slices
\[
 \{F(z,\cdot):z\in\mathbb D\}\subset H^\infty(\mathbb D)
 \tag{349.3}
\]
is totally bounded in the supremum norm. For a finite sum
$\sum_{\ell=1}^N a_\ell(z)b_\ell(w)$, these slices lie in a
bounded subset of a finite-dimensional space, hence are totally
bounded. Uniform approximation by such sums transfers that property
to (349.3): choose an approximant within half the desired covering
radius and a finite net for its slices.

Consider instead the bounded holomorphic function
\[
 F(z,w)=\exp\left(-\frac{1+(z+w)/2}{1-(z+w)/2}\right).
 \tag{349.4}
\]
The Cayley transform has positive real part on the disk, so
$|F|\leq1$. For each real $r<1$, $F(r,\cdot)$ extends holomorphically
past the closed disk: its only possible displayed singularity is at
$w=2-r>1$. As $r\uparrow1$, these slices converge locally uniformly to
\[
 F_1(w)=\exp\left(-\frac{3+w}{1-w}\right)
       =e^{-1}\exp\left(-2\frac{1+w}{1-w}\right).
\]
This limit has no continuous extension at $w=1$. Along the real
radius it tends to zero. Along the inverse Cayley images of
$1+it$, which remain inside the disk and tend to $1$, its modulus
is $e^{-3}$ and its phase oscillates. If (349.3) were totally
bounded, completeness of $H^\infty$ would give a uniformly convergent
subsequence of $F(r,\cdot)$. Its limit would be $F_1$ by local
convergence. But the disk algebra of functions continuous on the
closed disk is closed in this norm, a contradiction.

Thus (349.4) is outside $\mathcal S_2$. Regarding it as independent
of the remaining variables gives the same obstruction in every
$n>2$, by restricting the other variables to zero. The slice-algebra
theorem is therefore a proper subcase, not a disguised proof for all
of $H^\infty(\mathbb D^n)$.

\paragraph{Regularization solves individual problems but loses a bound.}
For $0<r<1$, set $f_{j,r}(z)=f_j(rz)$. These functions extend past
the closed polydisc and belong to $\mathcal S_n$, since their Taylor
polynomials converge uniformly there. Their common lower bound is
still $\delta$. The preceding theorem supplies bounded solutions
$g_{j,r}$, separately for every $r$.

If one could choose a sequence $r_\ell\uparrow1$ with
\[
 \sup_\ell\max_j\|g_{j,r_\ell}\|_\infty<\infty,
 \tag{349.5}
\]
Montel compactness and a diagonal subsequence over compact polydiscs
would give locally uniform limits $g_j\in A$. Passing to the finite
Bezout sum gives $\sum f_jg_j=1$. This proves that (349.5) suffices.
It is not supplied by qualitative solvability at each $r$.

There is a useful local stability statement. Given one solved tuple
$\sum f_jg_j=1$, perturb it to $f_j+h_j$ and suppose
$\|\sum h_jg_j\|_\infty<1$. Then
\[
 \widetilde g_j=\frac{g_j}{1+\sum_i h_ig_i}
 \tag{349.6}
\]
is bounded and solves the perturbed equation. The allowed perturbation
size depends on the already constructed solution norms. Radial
regularization is locally uniform but need not be uniform up to the
boundary, so (349.6) cannot repair the missing bound automatically.

\paragraph{Why solving on coordinate disks separately is insufficient.}
Fixing all but one coordinate and invoking disk corona gives a
solution on each slice. It does not select these solutions
holomorphically as functions of the remaining parameters. Taking
arbitrary choices may destroy even continuity. Taking continuous
partitions of unity makes the patched coefficients nonholomorphic;
their $\bar\partial$ errors then require a bounded correction with
the same Bezout constraints. This is precisely a new analytic
problem, not an algebraic consequence of separate solvability.

Likewise pointwise coefficients $\overline f_j/\sum|f_i|^2$ are
uniformly bounded and satisfy the Bezout equation, but are generally
not holomorphic. For the single datum $f(z)=2+z_1$, the reciprocal
is holomorphic whereas this distinction is hidden by the one-term
identity; with several independent data, the antiholomorphic
derivatives of the pointwise coefficients do not cancel individually.

\paragraph{Verification and first gap.}
The saved diagnostic samples (349.2) and (349.6) numerically for
finite complex inputs and checks the Cayley-path modulus in (349.4).
These are checks of the explicit formulas; total boundedness and the
maximal-ideal proof are analytic arguments above. The strongest
proved result is the corona theorem for $\mathcal S_n$, conditional
only on the classical disk theorem. The first missing general step
is a uniform regularized solution bound such as (349.5), or a bounded
holomorphic gluing construction outside that proper subalgebra.
The selected arbitrary-data assertion remains \textbf{UNRESOLVED}.
\subsection{Sources}
\begin{itemize}
\item[S1]Alexander Brudnyi and Mahishanka Withanachchi, Solvability of the Bézout Equation for Banach Algebra–valued H-infinity Functions on the Polydisk, arXiv2512.23733v2 (3Jan2026).
\url{https://arxiv.org/pdf/2512.23733v2}.
\textit{Formulation source.}
\end{itemize}
\end{document}
